\par\noindent\textbf{(1)} 写 \(|\psi(t)\rangle=\sum_n c_n(t)e^{-(i/\hbar)\int^t E_n(s)ds}|n(t)\rangle\)。代入 \(i\hbar|\dot\psi\rangle=H|\psi\rangle\) 得 \(\dot c_m=-\sum_nc_n\langle m|\dot n\rangle e^{(i/\hbar)\int^t(E_m-E_n)ds}\)。对角项产生几何相位，非对角项驱动跃迁。
\par\noindent\textbf{(2)} 对 \(H|n\rangle=E_n|n\rangle\) 求导并左乘 \(\langle m|\)（\(m\ne n\)），可得 \(\langle m|\dot n\rangle=\langle m|\dot H|n\rangle/(E_n-E_m)\)。因此跃迁振幅被瞬时能隙抑制，但随时间变化越快，\(\dot H\) 越大。
\par\noindent\textbf{(3)} 常用充分的局部慢变判据量纲为 \(\hbar|\langle m|\dot H|n\rangle|/|E_m-E_n|^2\ll1\)，并须兼顾总演化时间内的相位干涉和端点平滑性。它是估计而非所有时间依赖系统中严格充分的数学定理。
\par\noindent\textbf{(4)} 若能隙关闭，分母趋零，普通非简并绝热展开失效，需处理简并子空间或避免交叉；对闭合参数回路，沿单一本征态演化可积累Berry相位 \(\gamma_n=i\oint\langle n|\nabla_\lambda n\rangle\cdot d\boldsymbol\lambda\)，与动力相位不同且规范不变模 \(2\pi\)。
